\section{خلاصه‌ی کل کار: از مدل پایه تا \lr{v3}}
\label{sec:overview}

این فصل کل مسیر پروژه را در چند صفحه مرور می‌کند تا بدون خواندن همه‌ی فصل‌ها بتوان نتیجه را دید. پرسش همه‌ی مدل‌ها یکی است: آیا \lr{EEG} سطح سر چیزی درباره‌ی فعالیت عضله‌های دست می‌داند که از حرکت خود دست و از زمان‌بندی تکلیف نتوان فهمید؟ داده \lr{WAY-EEG-GAL} است: ۱۲ نفر که جسمی را بارها گرفته، بلند کرده و نگه داشته‌اند. هدف همه‌ی مدل‌ها پوش \lr{EMG} پنج عضله است و همه روی همان دو سری آزمون هر نفر سنجیده می‌شوند.

\subsection{مدل پایه (\lr{v1}): دقت خوب، سهم نامعلوم \lr{EEG}}

مدل پایه (\lr{KG-GT}) \lr{EEG} و سینماتیک را درون یک گراف با هم مخلوط می‌کرد و یک خروجی داشت. روی سری‌های آزمون به \lr{VAF} \lr{۷۲٫۲} درصد و $r$ برابر \lr{۰٫۸۵۱} رسید. آزمون‌های حذف ورودی نشان دادند تقریباً همه‌ی این دقت از سینماتیک می‌آید، و با یک خروجی نمی‌شد سهم \lr{EEG} را جدا اندازه گرفت (فصل~\ref{sec:gaps}).

\subsection{مدل \lr{v2}: سه خروجی جدا}

مدل \lr{v2} سه خروجی جدا ساخت تا این دو سهم از هم جدا شوند: \lr{K} فقط از سینماتیک و تاریخچه‌ی بلند کردن‌ها، \lr{E} فقط از \lr{EEG}، و \lr{F} که \lr{K} را با یک اصلاح کوچک از \lr{EEG} جمع می‌کند. نتیجه‌های اصلی:

\begin{itemize}
\item خروجی \lr{K} به \lr{VAF} \lr{۷۵٫۸} درصد و $r$ برابر \lr{۰٫۸۷۲} رسید و در هر ۱۲ نفر از مدل پایه بهتر بود.
\item خروجی \lr{E} به‌تنهایی به $r$ برابر \lr{۰٫۶۴۵} رسید و در هر ۱۲ نفر از آزمون \lr{null} گذشت.
\item خروجی \lr{F} با \lr{K} برابر ماند. \lr{gate} اصلاح \lr{EEG} از \lr{۰٫۰۴} بالاتر نرفت و \lr{EEG} فقط همبستگی حدود \lr{۰٫۰۷} با خطای \lr{K} داشت.
\end{itemize}

سه آزمایش کنترل (فصل~\ref{sec:controls}) این تصویر را دقیق‌تر کردند. یک مدل خطی که فقط زمان رخدادهای تکلیف را می‌داند (\lr{phase}) به $r$ برابر \lr{۰٫۸۱۴} می‌رسد و در هر ۱۲ نفر از \lr{E} بهتر است. فقط دانستن زمان روشن شدن \lr{LED} به \lr{۰٫۵۷۶} می‌رسد. پس بخش بزرگی از آنچه \lr{E} می‌داند، ساعت تکلیف است. در یک نسخه‌ی علّی که \lr{EMG} را چند صد میلی‌ثانیه جلوتر پیش‌بینی می‌کرد، سهم \lr{EEG} کمی بیشتر شد (\lr{F} منهای \lr{K} حدود \lr{+۰٫۵۵} واحد \lr{VAF} در افق ۵۰۰ میلی‌ثانیه).

\subsection{مدل \lr{v3}: \lr{EEG} پایه، سینماتیک فقط راهنمای آموزش}

یک هم‌تیمی در بازبینی \lr{v2} گفت جهت مدل برعکس است: قرار بود \lr{EEG} پایه باشد و سینماتیک فقط راهنمایی‌اش کند. \lr{v3} دو کار می‌کند (فصل~\ref{sec:v3}):

\begin{enumerate}
\item \textbf{مقایسه‌ی اصلی، پیشنهاد تیم (آزمایش بحرانی دوم).} برای هر نفر و هر افق پیش‌بینی (۰، ۲۰۰، ۵۰۰ و ۱۰۰۰ میلی‌ثانیه)، مدل \lr{B} از حرکت و نیروی گذشته به‌علاوه‌ی رخدادهای گذشته‌ی تکلیف ساخته می‌شود و مدل \lr{F} همان \lr{B} است به‌علاوه‌ی تاریخچه‌ی \lr{EEG}. دو مدل جز \lr{EEG} همه‌چیزشان یکی است. معیار اصلی $\text{\lr{eeg\_skill}} = 1 - \text{\lr{MSE}}(F) / \text{\lr{MSE}}(B)$ است: سهمی از خطای \lr{B} که \lr{EEG} برمی‌دارد. این بخش آموزش تازه لازم ندارد و از مدل‌های علّی \lr{v2} استفاده می‌کند.
\item \textbf{آزمایش ثانویه.} یک مدل فقط‌\lr{EEG} دو بار آموزش می‌بیند: بدون راهنما (\lr{E0}) و با جریان سینماتیک \lr{v2} به‌عنوان معلم ثابت که فقط در آموزش کمک می‌کند (\lr{Eguided}). در زمان آزمون این مدل فقط \lr{EEG} می‌بیند.
\end{enumerate}

برای اینکه ساعت تکلیف با اطلاعات واقعی \lr{EEG} اشتباه نشود، چند کنترل هم اضافه شد: مدل‌های زمان‌بندی علّی، جابه‌جایی \lr{EEG} بین \lr{trial}های هم‌شرط با هم‌تراز کردن همه‌ی رخدادها (\lr{phase-matched swap})، توان باند کانال‌های پیرامونی (چشم و گردن)، زمان شروع فعالیت عضله (\lr{onset}) و آزمون انحراف از الگوی میانگین \lr{trial}.

\subsection{نتیجه‌های تا امروز}

\begin{itemize}
\item \textbf{\lr{EEG} چیزی به \lr{B} اضافه نمی‌کند.} در هیچ‌کدام از چهار افق، \lr{F} از \lr{B} بهتر نیست. میانگین \lr{eeg\_skill} در هر چهار افق کمی منفی است: خطای \lr{F} بین \lr{۰٫۶} و \lr{۲٫۰} درصد بیشتر از \lr{B} است، و $p$ پس از تصحیح \lr{Holm} در همه‌ی افق‌ها برابر \lr{۱٫۰۰} است. مقدار \lr{eeg\_skill} فقط در ۳ تا ۶ نفر از ۱۲ نفر مثبت است. \lr{EEG} جابه‌جاشده (\lr{trial} دیگر با همان زمان‌بندی) و توان باند کانال‌های پیرامونی هم همان‌قدر نتیجه می‌دهند. در هیچ فاز تکلیفی هم اثر معنادار دیده نمی‌شود (فصل~\ref{sec:v3-bvsf}).
\item \textbf{رخدادهای گذشته کمی به حرکت اضافه می‌کنند.} در افق‌های ۲۰۰ و ۵۰۰ میلی‌ثانیه، \lr{B} خطای \lr{K} را حدود ۲ تا ۳ درصد کم می‌کند و زمان شروع فعالیت را بهتر پیش‌بینی می‌کند (در ۵۰۰ میلی‌ثانیه، میانه‌ی خطای \lr{onset} از ۴۰۰ به ۲۴۰ میلی‌ثانیه).
\item \textbf{راهنمای سینماتیک در \lr{validation} اثری ندارد.} هر ۱۴۴ \lr{run} آموزش (\lr{E0} و سه مقدار \lr{scale}، هر کدام با سه \lr{seed} برای ۱۲ نفر) تمام شده است. \lr{loss} \lr{validation} مدل‌های راهنماشده با \lr{E0} عملاً برابر است: بیشترین تغییر \lr{-۰٫۸۳} درصد با \lr{scale} \lr{۰٫۳} است، که در ۷ نفر از ۱۲ کمتر است و معنادار نیست ($p$ برابر \lr{۰٫۱۰}).
\item \textbf{هنوز اجرا نشده:} ارزیابی \lr{E0} و \lr{Eguided} روی سری‌های آزمون (گام ۷) در حال اجراست، و حکم نهایی پرسش ثانویه بعد از آن می‌آید (فصل~\ref{sec:v3-pending}).
\end{itemize}

\paragraph{در یک جمله.} با این داده و این روش، وقتی حرکت گذشته و زمان رخدادهای تکلیف را داریم، تاریخچه‌ی \lr{EEG} در هیچ افقی از صفر تا یک ثانیه چیزی به پیش‌بینی فعالیت عضله اضافه نمی‌کند. سهم کوچکی که \lr{v2} در افق‌های بلند دیده بود، با افزودن رخدادهای گذشته به مبنا از بین می‌رود. این با این فرض سازگار است که آن سهم از دانستن زمان تکلیف می‌آمد.

\subsection{پایلوت خطای پیش‌بینی و نتیجه‌گیری نهایی}

پس از نتیجه‌ی منفی \lr{v3}، یک پرسش دیگر از همین داده آزموده شد: وقتی جسم سنگین‌تر یا سبک‌تر از انتظار است، \lr{EEG} چه نشان می‌دهد (فصل~\ref{sec:pe-pilot})؟
\begin{itemize}
\item رفتار الگوی کلاسیک را در هر ۱۲ نفر نشان می‌دهد.
\item \lr{EEG} به جسم سبک‌تر از انتظار پاسخ می‌دهد، ولی بعد از اینکه نیرو و \lr{EMG} فرق کرده‌اند.
\item وزن مورد انتظار پیش از تماس از \lr{EEG} رمزگشایی نمی‌شود.
\end{itemize}
فصل~\ref{sec:final} همه‌ی نتیجه‌ها را جمع می‌کند و مقاله‌ی ممیزی را به‌عنوان راه پیشنهادی توضیح می‌دهد.

\newpage
